% GENERATED FILE. Edit canonical lessons, then run npm run generate:books.
% canonical-source-hash: fnv1a64:bed8f29a4bc07667
% canonical-lessons: TA-C337-cati, TA-C337-culcci, TA-C337-kattikkotu, TA-C337-vanci, TA-C337-maruvetam

\chapter{Plots and Betrayal}
\label{ch:ta-plots-and-betrayal}
\hlchaptermodality{337}

\begin{chapteropening}
\textbf{By the end of this chapter:} \emph{I can tell a story of plots and betrayal.}

Complete the last lesson of chapter 337: I can tell a story of plots and betrayal.
\end{chapteropening}

\section[\texorpdfstring{\ta{சதி} (cati)}{சதி (cati)}]{\ta{சதி} (cati) — a conspiracy; a plot}

\label{lesson:TA-C337-cati}



\begin{quote}
Before the new one: say the Tamil for to submit, then the Tamil for to occupy.
\end{quote}

\subsection*{You'll want to know: \ta{சதி}}
\textbf{\ta{சதி}} — \emph{cati} — \textquotedblleft{}a conspiracy; a plot\textquotedblright{}.

\begin{quote}
\ta{மன்னரைக்} \ta{கொல்ல} \ta{ஒரு} \ta{சதி} \ta{நடந்தது}.
\end{quote}

\begin{quote}
\emph{maṉṉaraik kolla oru cati naṭantatu.}
\end{quote}

\begin{quote}
There was a plot to kill the king.
\end{quote}

\ta{சதி} \ta{செய்}, to conspire.

\begin{cousinweb}
\textbf{\ta{சதி}} (\emph{cati}) is native Dravidian (DEDR 2323): to deceive, and as a noun treachery, wiles.
\end{cousinweb}

\begin{grammarlens}[title={What you've built}]
One more word for telling what happened.
\end{grammarlens}

\subsection*{Your turn}
\begin{itemize}
\raggedright
  \item \emph{Say it:} \emph{cati}
  \item \emph{Say it:} \emph{cati}, once more
  \item \emph{Say it:} \emph{maṉṉaraik kolla oru cati naṭantatu.} — There was a plot to kill the king.
  \item \emph{Recall:} say the Tamil for to submit, then the Tamil for to occupy, then say \emph{cati} again
\end{itemize}

\begin{tcolorbox}[breakable,colback=teal!4,colframe=teal!35!black,title={Before you move on}]
What does \ta{சதி} mean? (\textquotedblleft{}A conspiracy; a plot\textquotedblright{}.) Say it once more.
\end{tcolorbox}
\section[\texorpdfstring{\ta{சூழ்ச்சி} (cūḻcci)}{சூழ்ச்சி (cūḻcci)}]{\ta{சூழ்ச்சி} (cūḻcci) — a scheme; a ploy}

\label{lesson:TA-C337-culcci}



\begin{quote}
Before the new one: say the Tamil for to occupy, then the Tamil for a conspiracy.
\end{quote}

\subsection*{You'll want to know: \ta{சூழ்ச்சி}}
\textbf{\ta{சூழ்ச்சி}} — \emph{cūḻcci} — \textquotedblleft{}a scheme; a ploy\textquotedblright{}.

\begin{quote}
\ta{தளபதியின்} \ta{சூழ்ச்சி} \ta{கடைசியில்} \ta{வெளிப்பட்டது}.
\end{quote}

\begin{quote}
\emph{taḷapatiyiṉ cūḻcci kaṭaiciyil veḷippaṭṭatu.}
\end{quote}

\begin{quote}
The commander's scheme came out in the end.
\end{quote}

\ta{சூழ்ச்சி} \ta{செய்}, to scheme.

\begin{cousinweb}
\textbf{\ta{சூழ்ச்சி}} (\emph{cūḻcci}) is native Dravidian: DEDR 2735 gives it under \ta{சூழ்}, \textquotedblleft{}to deliberate, conspire\textquotedblright{}, as \textquotedblleft{}deliberation, counsel\textquotedblright{}.
\end{cousinweb}

\begin{grammarlens}[title={What you've built}]
One more word for telling what happened.
\end{grammarlens}

\subsection*{Your turn}
\begin{itemize}
\raggedright
  \item \emph{Say it:} \emph{cūḻcci}
  \item \emph{Say it:} \emph{cūḻcci}, once more
  \item \emph{Say it:} \emph{taḷapatiyiṉ cūḻcci kaṭaiciyil veḷippaṭṭatu.} — The commander's scheme came out in the end.
  \item \emph{Recall:} say the Tamil for to occupy, then the Tamil for a conspiracy, then say \emph{cūḻcci} again
\end{itemize}

\begin{tcolorbox}[breakable,colback=teal!4,colframe=teal!35!black,title={Before you move on}]
What does \ta{சூழ்ச்சி} mean? (\textquotedblleft{}A scheme; a ploy\textquotedblright{}.) Say it once more.
\end{tcolorbox}
\section[\texorpdfstring{\ta{காட்டிக்கொடு} (kāṭṭikkoṭu)}{காட்டிக்கொடு (kāṭṭikkoṭu)}]{\ta{காட்டிக்கொடு} (kāṭṭikkoṭu) — to betray}

\label{lesson:TA-C337-kattikkotu}



\begin{quote}
Before the new one: say the Tamil for a conspiracy, then the Tamil for a scheme.
\end{quote}

\subsection*{You'll want to know: \ta{காட்டிக்கொடு}}
\textbf{\ta{காட்டிக்கொடு}} — \emph{kāṭṭikkoṭu} — \textquotedblleft{}to betray\textquotedblright{}.

\begin{quote}
\ta{ஒரு} \ta{நண்பனே} \ta{அவரைக்} \ta{காட்டிக்கொடுத்தான்}.
\end{quote}

\begin{quote}
\emph{oru naṇpaṉē avaraik kāṭṭikkoṭuttāṉ.}
\end{quote}

\begin{quote}
It was a friend who betrayed him.
\end{quote}

Past (I …): \ta{காட்டிக்கொடுத்தேன்}, \emph{kāṭṭikkoṭuttēṉ}.

Literally to show and give: to give someone away.

\begin{cousinweb}
\textbf{\ta{காட்டிக்கொடு}} (\emph{kāṭṭikkoṭu}) is a Tamil compound built from \ta{காட்டு}, \textquotedblleft{}to show\textquotedblright{} (DEDR 1443), and \ta{கொடு}, \textquotedblleft{}to give\textquotedblright{} (DEDR 2053).
\end{cousinweb}

\begin{grammarlens}[title={What you've built}]
One more word for telling what happened.
\end{grammarlens}

\subsection*{Your turn}
\begin{itemize}
\raggedright
  \item \emph{Say it:} \emph{kāṭṭikkoṭu}
  \item \emph{Say it:} \emph{kāṭṭikkoṭu}, once more
  \item \emph{Say it:} \emph{oru naṇpaṉē avaraik kāṭṭikkoṭuttāṉ.} — It was a friend who betrayed him.
  \item \emph{Recall:} say the Tamil for a conspiracy, then the Tamil for a scheme, then say \emph{kāṭṭikkoṭu} again
\end{itemize}

\begin{tcolorbox}[breakable,colback=teal!4,colframe=teal!35!black,title={Before you move on}]
What does \ta{காட்டிக்கொடு} mean? (\textquotedblleft{}To betray\textquotedblright{}.) Say it once more.
\end{tcolorbox}
\section[\texorpdfstring{\ta{வஞ்சி} (vañci)}{வஞ்சி (vañci)}]{\ta{வஞ்சி} (vañci) — to deceive}

\label{lesson:TA-C337-vanci}



\begin{quote}
Before the new one: say the Tamil for a scheme, then the Tamil for to betray.
\end{quote}

\subsection*{You'll want to know: \ta{வஞ்சி}}
\textbf{\ta{வஞ்சி}} — \emph{vañci} — \textquotedblleft{}to deceive\textquotedblright{}.

\begin{quote}
\ta{அந்த} \ta{மந்திரி} \ta{மன்னரை} \ta{வஞ்சித்தார்}.
\end{quote}

\begin{quote}
\emph{anta mantiri maṉṉarai vañcittār.}
\end{quote}

\begin{quote}
That minister deceived the king.
\end{quote}

Past (I …): \ta{வஞ்சித்தேன்}, \emph{vañcittēṉ}.

A written word; the noun is \ta{வஞ்சனை}, deceit.

\begin{cousinweb}
\textbf{\ta{வஞ்சி}} (\emph{vañci}) is a Tamil verb built on Sanskrit \emph{vañc}, \textquotedblleft{}to cheat, deceive\textquotedblright{} (Monier-Williams).
\end{cousinweb}

\begin{grammarlens}[title={What you've built}]
One more word for telling what happened.
\end{grammarlens}

\subsection*{Your turn}
\begin{itemize}
\raggedright
  \item \emph{Say it:} \emph{vañci}
  \item \emph{Say it:} \emph{vañci}, once more
  \item \emph{Say it:} \emph{anta mantiri maṉṉarai vañcittār.} — That minister deceived the king.
  \item \emph{Recall:} say the Tamil for a scheme, then the Tamil for to betray, then say \emph{vañci} again
\end{itemize}

\begin{tcolorbox}[breakable,colback=teal!4,colframe=teal!35!black,title={Before you move on}]
What does \ta{வஞ்சி} mean? (\textquotedblleft{}To deceive\textquotedblright{}.) Say it once more.
\end{tcolorbox}
\section[\texorpdfstring{\ta{மாறுவேடம்} (māṟuvēṭam)}{மாறுவேடம் (māṟuvēṭam)}]{\ta{மாறுவேடம்} (māṟuvēṭam) — a disguise}

\label{lesson:TA-C337-maruvetam}



\begin{quote}
Before the new one: say the Tamil for to betray, then the Tamil for to deceive.
\end{quote}

\subsection*{You'll want to know: \ta{மாறுவேடம்}}
\textbf{\ta{மாறுவேடம்}} — \emph{māṟuvēṭam} — \textquotedblleft{}a disguise\textquotedblright{}.

\begin{quote}
\ta{மன்னர்} \ta{மாறுவேடத்தில்} \ta{நகரத்தைச்} \ta{சுற்றினார்}.
\end{quote}

\begin{quote}
\emph{maṉṉar māṟuvēṭattil nakarattaic cuṟṟiṉār.}
\end{quote}

\begin{quote}
The king went round the city in disguise.
\end{quote}

\ta{மாறுவேடம்} \ta{போடு}, to put on a disguise.

\begin{cousinweb}
\textbf{\ta{மாறுவேடம்}} (\emph{māṟuvēṭam}) joins native \ta{மாறு}, \textquotedblleft{}to be changed\textquotedblright{} (DEDR 4834), with \ta{வேடம்}, \textquotedblleft{}costume\textquotedblright{}, borrowed from Sanskrit \emph{vēṣa}, \textquotedblleft{}dress, assumed appearance\textquotedblright{} (CDIAL 12129).
\end{cousinweb}

\begin{grammarlens}[title={What you've built}]
One more word for telling what happened.
\end{grammarlens}

\subsection*{Your turn}
\begin{itemize}
\raggedright
  \item \emph{Say it:} \emph{māṟuvēṭam}
  \item \emph{Say it:} \emph{māṟuvēṭam}, once more
  \item \emph{Say it:} \emph{maṉṉar māṟuvēṭattil nakarattaic cuṟṟiṉār.} — The king went round the city in disguise.
  \item \emph{Recall:} say the Tamil for to betray, then the Tamil for to deceive, then say \emph{māṟuvēṭam} again
\end{itemize}

\begin{tcolorbox}[breakable,colback=teal!4,colframe=teal!35!black,title={Before you move on}]
What does \ta{மாறுவேடம்} mean? (\textquotedblleft{}A disguise\textquotedblright{}.) Say it once more.
\end{tcolorbox}
